% RESEARCH_DRAFT_SHA256: f7526526a289eff746bb483325f731461f28f72b54505cb5f1f820a27fbfe2cc
% Layout and source-reference wrappers expanded for the web renderer.
\documentclass[11pt]{article}
\usepackage[margin=25mm]{geometry}
\usepackage{fontspec,unicode-math}
\setmainfont{Latin Modern Roman}
\setmathfont{Latin Modern Math}
\usepackage{amsmath,amssymb,mathtools}
\usepackage{xurl,hyperref}
\hypersetup{colorlinks=true,urlcolor=blue,linkcolor=blue}
\setcounter{secnumdepth}{0}
\setlength{\parindent}{0pt}
\setlength{\parskip}{5pt}
\setlength{\emergencystretch}{3em}
\newcommand{\sref}[2]{\hyperlink{source:#1:#2}{[S#2]}}
\newcommand{\cataloguescope}[1]{\paragraph{Source question.}#1}
\newcommand{\floor}[1]{\left\lfloor#1\right\rfloor}
\begin{document}
\section*{Research-attempt provenance}
This standalone research attempt retains the complete original definition body below,
including its historical source and preparation notes. Those notes describe the
earlier catalogue preparation; the new attempt and its verification are identified
separately after that body. No original ID, catalogue status, or source statement is changed.

\section{Computing Hecke Operators for Arithmetic-Group Cohomology}\label{statement}

\subsection{Definitions and mathematical statement}

\paragraph{Shared notation.}Standard number systems are denoted by $\mathbb N,\mathbb Z,\mathbb Q,\mathbb R,\mathbb C$; index ranges, computational models and specialized notation follow the cited problem source and the definitions below.


Arithmetic-group cohomology is the cohomology of the specified arithmetic group or associated quotient. A Hecke operator is a correspondence-induced linear action on it; eigenvalues are its scalar actions on eigenvectors.
\paragraph{Statement recorded in the supplied source.}
Can Hecke operators and their eigenvalues be computed for all arithmetic groups for which the cohomology can be calculated? \href{https://doi.org/10.4171/owr/2016/52}{[S2]}
\paragraph{Scope qualification.}
The group family, cohomology coefficients, degree, and input representation must be read from the source.
\paragraph{Source status.}
The supplied catalogue labels this record partially solved. Its source-status metadata was checked against the supplied snapshot on 2026-10-11. This import does not independently establish current open status or mathematical correctness. The catalogue literature assessment is dated 2026-08-21.
\subsection{Short English statement}
Compute Hecke operators and eigenvalues whenever the underlying arithmetic-group cohomology is computable.
\subsection{Sources}
\begin{itemize}
\item[S1] Ulam AI and the UnsolvedMath Contributors, UnsolvedMath: A Curated Collection of Open Mathematics Problems, catalogue entry. Catalogue curation and metadata: CC BY 4.0. The supplied source question is retained; definitions and the English summary are editorial additions.
\url{https://huggingface.co/datasets/ulamai/UnsolvedMath}.
\item[S2] Mini-Workshop: Computations in the Cohomology of Arithmetic Groups (2017), Oberwolfach Reports, DOI 10.4171/owr/2016/52, pp. 2941--2973
\url{https://doi.org/10.4171/owr/2016/52}.
\item[S3] McConnell and MacPherson
\url{https://arxiv.org/abs/2010.06036}.
\textit{Underlying Oberwolfach Reports material retains its source-specific licence, including CC BY-SA 4.0 where applicable; see THIRD\_PARTY\_NOTICES.md.}
\end{itemize}
\paragraph{Preparation and review.}
Original question attributed to the cited source and UnsolvedMath contributors. Codex (GPT-6; exact serving revision not exposed) assisted with definitions, English summaries, source inspection and formatting. No new human mathematical authorship or independent proof verification is asserted. Catalogue identity checked against the complete supplied registry and existing repository statements; known flagged repetitions and obvious identical conjectures excluded. Full mathematical-equivalence/literature review has not been performed. Statement preparation and final review followed contribution prompts C and D; automated checks verify structure and identity, not open status.
\label{end:statement}

\section{Literature and context check}
The primary problem session in OWR 52/2016, pp.~2970--2971, was read. Harder's question explicitly seeks implementations extending cohomology software. The abstract of McConnell--MacPherson, arXiv:2010.06036v2, was also read: it reports an algorithm for arithmetic subgroups of general linear groups, finite-dimensional local coefficients, and all cohomological degrees, using a well-tempered complex. Its construction and proof were not rederived here. The calculation below isolates what ordinary linear algebra can do once the correspondence has already been represented; it supplies no replacement for that construction.
\section{Research attempt: Cohomology action from supplied chain maps}

\paragraph{Conditional exact algorithm.}
Let a finite cochain complex over $\mathbb Q$ be supplied by rational matrices $d^j:C^j\to C^{j+1}$ with $d^{j+1}d^j=0$. Suppose rational matrices $T^j:C^j\to C^j$ are also supplied and satisfy
\[
 d^jT^j=T^{j+1}d^j.
\]
Find by elimination a basis matrix $B$ for $\operatorname{im}d^{j-1}$ and extend it inside $\ker d^j$ to a basis $[B\ H]$. Because $T^j$ preserves cycles and boundaries, there are unique matrices $U,M$ such that
\[
 T^jH=BU+HM.
\]
The induced operator on $H^j=\ker d^j/\operatorname{im}d^{j-1}$ is exactly $M$: each column equation gives the image of one chosen cohomology class. Gaussian elimination constructs all of these rational matrices in finitely many arithmetic operations; this is not a uniform bit-complexity bound for constructing the input complex. The characteristic polynomial $\det(zI-M)$ is an exact finite representation of the multiset of algebraic eigenvalues. Its roots, with multiplicity, specify those eigenvalues without pretending that arbitrary algebraic numbers have rational closed forms. The zero-dimensional cohomology case has the empty matrix and characteristic polynomial $1$.

\paragraph{A computed cancellation test.}
Even when a complex can be compressed, deleting the cancelled coordinates of $T$ is wrong in general. Use a chain complex $C_1=\mathbb Q$, $C_0=\mathbb Q^2$, $d(a)=(a,a)$, and chain maps
\[
 T_1=(2),\qquad T_0=\begin{pmatrix}1&1\\0&2\end{pmatrix}.
\]
Indeed $T_0d=dT_1$. The quotient $H_0=C_0/\operatorname{im}d$ is represented by $e_2$, with projection $p(x,y)=y-x$ and inclusion $i(z)=(0,z)$. Its induced action is
\[
 pT_0i(z)=z.
\]
Simply keeping the bottom-right entry of $T_0$ instead gives $2$, a wrong eigenvalue. The missing $-x$ in the projection records the boundary relation $e_1=-e_2$.

\paragraph{What cohomology alone fails to determine.}
Consider the one-term complex $C^0=\mathbb Q$ with zero differential. For every rational $\lambda$, the map $T^0=\lambda I$ is a chain map and induces eigenvalue $\lambda$ on the same one-dimensional cohomology. Thus the vector space, or even this entire differential complex, contains no information selecting a particular correspondence operator. This example does not assert that arbitrary $\lambda$ occurs as an arithmetic Hecke operator; it pinpoints why a procedure receiving only cohomology data cannot recover the missing correspondence. Actual double-coset representatives, transfer maps, coefficient actions, and compatible chain-level reductions must be constructed and verified before the conditional algorithm applies.
\section{Verification and adversarial audit}
The cycle/boundary invariance is a direct consequence of the displayed chain-map equation. The two-by-two example was multiplied explicitly and the quotient projection checked on the boundary $(1,1)$. The one-term example concerns insufficient input data, not noncomputability of arithmetic Hecke actions. Exact algebraic eigenvalue representation is distinguished from rational eigenvalues or a floating-point eigensolver.
\paragraph{Audit corrections.} Replaced an unjustified inference from computable cohomology to computable Hecke maps by a precise supplied-chain-map hypothesis. Added an explicit example where naive coordinate deletion changes an eigenvalue.
\paragraph{External computation.} None; symbolic arguments and boundary checks only.
\section{Scope and remaining gaps}
The first missing step is an effective chain-level realization of every desired Hecke correspondence for every original arithmetic group and coefficient system. Supplying that step as an input is a strict extra hypothesis. The later primary abstract covers an important named family, but neither it nor the conditional linear algebra is claimed to settle all groups in the source question.
\paragraph{Integration review.} Codex (GPT-6; exact serving revision not exposed) reviewed the written arguments and expanded layout and source-reference wrappers. Repository metadata and the production web renderer were checked using Python 3.13 and JavaScript with MathJax 3.2.2. This model-assisted review is not independent human proof verification; no PDF compilation or formal proof checking is asserted.

\section{Completion Estimate}
\noindent\textbf{COMPLETION: 0\%}
The calculations below the preserved source statement investigate a concrete route,
but establish no substantiated new progress toward the entire original target.
Known special cases and failed sufficient conditions are not credited as a solution.
\section{References}
The original source references and catalogue attribution are preserved in the
Sources section above. The following URLs were actually consulted for this attempt;
retention of an additional original reference does not claim fresh consultation of it.
\begin{itemize}
\item \url{https://ems.press/content/serial-article-files/46655} (consulted 2026-10-11). 
\item \url{https://arxiv.org/abs/2010.06036} (consulted 2026-10-11). 
\end{itemize}
\section{Model and contributor disclosure}
OpenAI gpt-6-astra, requested reasoning configuration ultra, generated this research
attempt on 2026-10-11. The exact serving revision is not exposed. Codex GPT-6
performed orchestration, repository-format review, and written-proof review of this attempt; its exact serving revision
is also not exposed. The mathematical route, boundary checks and final adversarial
self-audit followed contribution prompts A then B in
\texttt{docs/contribution\_prompts.md}, after reading \texttt{CONTRIBUTING.md}.
Python 3.14.2 assembled the standalone documents; any mathematical computation is
specified in Verification. No theorem prover was used. No human mathematical
authorship, substantial human proof correction, or independent human proof checking
is asserted. Model-assisted review, source retrieval, and structural checks do not
establish correctness or novelty. Full private reasoning transcripts are unavailable
as reproducibility artifacts; the complete arguments, inputs and outcomes actually
used are stated here or in the companion batch audit.
\section{Final status}
Unresolved as a new research contribution. This is a concrete route investigation,
not a claim of new complete resolution or a current-status certification. A known
answer reproduced for context does not become a new contribution. The catalogue's
source-status label is preserved; this attempt status does not assert that a
source-labelled solved question is open.
\end{document}
